\documentclass[11pt]{article}
\usepackage[margin=0.82in]{geometry}
\usepackage{amsmath}
\usepackage{booktabs}
\usepackage{array}
\usepackage{float}
\newcolumntype{L}[1]{>{\raggedright\arraybackslash}p{#1}}

\title{Project Three: Statistical Integrity Tests for Fund X}
\author{Prepared by Harman Singh for Professor Phelim Boyle}
\date{June 2026}

\begin{document}
\maketitle

\section*{Data Box}
The analysis uses 113 monthly observations from January 2017 through May 2026. Returns are measured in percent per month.

\begin{table}[H]
\centering
\caption{Summary Statistics Used Across Tests}
\begin{tabular}{lrr}
\toprule
Statistic & Fund X & S\&P 500 TR \\
\midrule
Sample mean (\%) & 3.50 & 1.33 \\
Sample standard deviation (\%) & 4.01 & 4.54 \\
Skewness & 2.65 & -0.45 \\
Excess kurtosis & 8.52 & 0.52 \\
Positive months & 108 & 80 \\
Negative months & 5 & 33 \\
Monthly Sharpe ratio & 0.87 & 0.29 \\
\bottomrule
\end{tabular}
\end{table}

\section*{Test 1: Wald-Wolfowitz Runs Test}
Fund X has 11 observed sign runs, with 108 positive months and 5 negative months.
The expected number of runs under independent signs is 10.56, with variance 0.73.
The maximum possible number of runs is $R_{max}=2n_-+1=11$, and Fund X reaches that maximum.

Because $n_-=5$ is very small, the normal approximation is not the decision rule.
The exact hypergeometric runs distribution has support $R\in\{2,\ldots,11\}$.
For the smoothing alternative, the relevant one-sided concern is too few runs, so the exact p-value is $P(R\leq 11)=1.000$.
The opposite upper-tail probability is $P(R\geq 11)=0.757$.

The exact test does not flag Fund X. In fact, the five negative months are as separated as possible.
The key nuance is power: once the analysis conditions on only five negative months, a runs test can judge whether losses cluster, but it cannot explain why the number of negative months is so small.

\clearpage
\section*{Test 2: Bollen-Pool Discontinuity Test}
Assuming normal returns with Fund X's mean and standard deviation, the probability of a negative month is 0.192.
This implies 21.66 expected negative months, compared with 5 observed negative months.
The binomial lower-tail p-value is 4.64e-06.
The S\&P 500 benchmark has $p_-=0.385$, 43.51 expected negative months, 33 observed negative months, and p-value 0.025.

Fund X's near-zero count is 5 months in $(0\%,0.5\%)$ and 3 months in $(-0.5\%,0\%)$, for a near-zero ratio of 1.67.
The five negative months are 2017-06 (-0.35\%); 2017-08 (-1.26\%); 2022-06 (-0.32\%); 2022-09 (-0.21\%); 2024-05 (-2.20\%).
Treating the near-zero window as the direct discontinuity check, there are 5 small positive months and 3 small negative months within $\pm 0.5\%$.
The exact conditional p-value for too many small positives is 0.363, and the two-sided exact p-value is 0.727.

The normal-count version is therefore best read as a stress diagnostic, not as final evidence of manipulation.
Fund X has skewness 2.65 and excess kurtosis 8.52, so normality is a poor approximation.
A 3-month circular block bootstrap, which preserves some dependence and the empirical skewed return shape, gives a median negative count of 5.00 and a 90\% interval of [2.00, 9.00].
The observed count of 5 negative months is at empirical lower-tail probability 0.615 under this calibration.

\section*{Test 3: Serial Autocorrelation Test}
The first three raw autocorrelations for Fund X are $\rho_1=0.347$, $\rho_2=0.157$, and $\rho_3=0.222$.
The first-order t-statistic is 3.90, with p-value 0.0002.
The S\&P 500 first-order autocorrelation is -0.138.

The active return series, Fund X minus S\&P 500 TR, has autocorrelations $\rho_1=0.013$, $\rho_2=-0.026$, and $\rho_3=0.080$.
Its first-order t-statistic is 0.14, with p-value 0.891.
The Ljung-Box p-value through 12 lags is 0.0001 for raw Fund X returns and 0.857 for active returns.

The single-lag t-test is an approximate test of $H_0:\rho_1=0$; a small p-value means the observed first-order autocorrelation would be unusual under no serial dependence.
The Ljung-Box statistic is broader because it tests whether the first 12 autocorrelations are jointly zero.
This pattern supports level-persistence rather than return smoothing. The raw Fund X returns are serially correlated, but the active returns are not.
The naive annualised Sharpe ratio is 3.02.
Using Lo's 12-lag adjustment gives 1.76; using 6 lags gives 2.05.
Both adjusted values remain economically exceptional.

\section*{Test 4: Abdulali Bias Ratio}
Using Fund X's sample standard deviation of 4.01\%, the upper interval $[0,\hat\sigma]$ contains 80 returns and the lower interval $[-\hat\sigma,0)$ contains 5 returns.
The observed bias ratio is therefore 13.33.
Under a normal distribution with Fund X's monthly Sharpe ratio, the expected bias ratio is approximately 2.11, already above Abdulali's usual warning thresholds for low-Sharpe hedge funds.

The bootstrap calibration uses 100,000 resamples of the 113 observed Fund X returns.
The bootstrap median is 13.17, the mean is 15.23, and the 90\% interval is [7.30, 28.67].
The observed bias ratio is at the 51.9th percentile of the bootstrap distribution.
The S\&P 500 bias ratio is 2.85.
The high Fund X bias ratio is typical for its own return distribution and reflects genuine high positive-return frequency, not a small-positive bunching anomaly.

\clearpage
\section*{Overall Assessment}
\begin{table}[H]
\centering
\caption{Summary of Statistical Integrity Tests}
\begin{tabular}{L{0.19\linewidth}L{0.24\linewidth}L{0.25\linewidth}L{0.21\linewidth}}
\toprule
Test & Key statistic & Assessment & Caveat \\
\midrule
Runs test & $R=11$, $R_{max}=11$ & No evidence of too few runs & Few negative months limits test power \\
Bollen-Pool & Normal count p=4.64e-06; local p=0.727 & No local discontinuity after exact near-zero check & Normal count test is assumption-sensitive \\
Serial autocorrelation & Raw $\rho_1=0.347$; active $\rho_1=0.013$ & Supports level-persistence, not smoothing & p-values are approximate dependence diagnostics \\
Bias ratio & BR=13.33; bootstrap percentile 51.9 & Typical for Fund X's own distribution & Abdulali thresholds are not calibrated to high-Sharpe funds \\
\bottomrule
\end{tabular}
\end{table}

Taken together, the tests do not provide statistical evidence of return smoothing or mechanical avoidance of losses.
They cannot rule out every possible concern about Fund X, but the more assumption-aware versions of the tests do not support the specific irregularities these diagnostics were designed to detect.

\end{document}
